\section{Complete recursive decomposition of a primitive projective partner}
\label{sec:recursive-primitive-partner}

\paragraph{Status and provenance.}
D-TD11--13, reserved question Q2dv, supplies the completeness interface
which was deliberately a hypothesis in HD-FT20--22.  The parent proposed
retaining and recursively analyzing every multiplicity action.  The full
characteristic-zero argument below is written by
\texttt{/root/imprimitive}, using the previously checked homogeneous
restriction and matrix-algebra arguments of D-BI19 and D-TD04.
The author, root, and independent reviewer have each read the complete proof and source interfaces; all whole-proof checks pass without repair. No numerical certificate was needed or executed.
No assertion that the
generalized Fitting subgroup acts irreducibly is used.

All vector spaces are finite dimensional over $\mathbf C$.  A finite
subgroup $P\subset\operatorname{PGL}(V)$ is called primitive if a linear
lift acts irreducibly and preserves no direct-sum decomposition with
at least two nonzero summands which it permutes.  These properties do
not depend on scalar choices.  For a nonscalar finite-projective-order
operator $T$, write
\[
 \nu_V(T)=\dim V-\max_{z\in\mathbf C^*}\dim\ker(T-zI).
 \tag{TDR.1}
\]
A represented projective Heisenberg group means an elementary abelian
subgroup $Q\subset\operatorname{PGL}(W)$ with a nondegenerate scalar
commutator pairing, whose linear inverse image acts irreducibly on $W$.
No splitting of its projective normalizer is included in this definition.

\subsection*{D-TD11. The multiplicity quotient is finite and primitive}

\begin{lemma}[D-TD11]
Let $P\subset\operatorname{PGL}(V)$ be finite and primitive, with
$d=\dim V>1$.  Let $N$ be any nontrivial normal subgroup of $P$.
There are vector spaces $W,U$, an isomorphism $V=W\otimes U$, and
genuine projective homomorphisms
\[
 \pi_W:P\longrightarrow\operatorname{PGL}(W),\qquad
 \pi_U:P\longrightarrow\operatorname{PGL}(U)
 \tag{TDR.2}
\]
with finite images, such that the following hold.
\begin{enumerate}
\item The inverse image of $N$ in the finite special-linear lift of
$P$ acts as a faithful irreducible group on $W$ and trivially on $U$.
In particular $\dim W\ge2$ and $\dim U\le d/2$.
\item Both projective images in (TDR.2) are primitive when their
spaces have dimension greater than one.  If $\dim U>1$, its image
is nontrivial, so it is an eligible strictly smaller recursion input.
\item Every original operator is a tensor product of representatives
of (TDR.2), up to its original scalar.  No coherent linear choice
of either factor is required.
\end{enumerate}
The same primitivity conclusion applies to either factor of any
factorwise tensor decomposition of a finite primitive coefficient
representation.
\end{lemma}

\begin{proof}
The inverse image $J$ of $P$ under
$\operatorname{SL}(V)\to\operatorname{PGL}(V)$ is finite, with
central scalar kernel $Z=\mu_d$.  Let $R$ be the inverse image of
$N$.  Thus $R\triangleleft J$, $Z\subset R$, and $R/Z=N$.
The canonical isotypic decomposition for $R$ is permuted by $J$;
irreducibility makes that permutation transitive.  Primitivity
forces a single isotypic summand.  This is exactly the homogeneous
alternative in Serre's normal-subgroup argument
\cite[Section~8.1, Proposition~24 and proof, p.~61]{OG:Serre1977};
its canonical summands are those of
\cite[Section~2.6, Theorem~8, p.~21]{OG:Serre1977}.

Choose the irreducible $R$-module $W$ in this restriction and let
$U=\operatorname{Hom}_R(W,V)$.  Evaluation gives
\[
 V\simeq W\otimes U,\qquad
 x|_V=\rho_W(x)\otimes I_U\quad(x\in R).
 \tag{TDR.3}
\]
The representation $\rho_W$ is faithful: an element in its kernel
acts trivially on all of $V$, and $J$ is an actual matrix group.
If $\dim W=1$, every element of $R$ acts by a scalar on $V$.
Then $R\subset Z$, contradicting $N\ne1$.  Hence $\dim W\ge2$.

By the irreducible matrix-algebra theorem
\cite[Section~6.2, Proposition~10 and proof, p.~48]{OG:Serre1977},
the algebra generated by $R$ on $V$ is
\[
 \mathcal A=\operatorname{End}(W)\otimes I_U,
 \qquad
 \operatorname{Cent}_{\operatorname{End}(V)}(\mathcal A)
       =I_W\otimes\operatorname{End}(U).
 \tag{TDR.4}
\]
Normality of $R$ makes $J$ normalize $\mathcal A$.  Every complex
matrix-algebra automorphism is inner: if $f_{ij}$ are the images
of the standard matrix units, choose a nonzero vector $v$ in
$\operatorname{im}f_{11}$; the vectors $f_{i1}v$ form a basis
in which the $f_{ij}$ are the standard matrix units.  Applying this
to the action of $g\in J$, and then using the centralizer in
(TDR.4), gives
\[
                         g=A_g\otimes B_g.
 \tag{TDR.5}
\]
The projective factors are unique.  Indeed the intersection of
the two algebras in (TDR.4) consists of scalar operators, so two
such expressions differ only by inverse scalar changes in their
two factors.  Multiplication therefore makes
$g\mapsto[A_g]$ and $g\mapsto[B_g]$ homomorphisms.  They kill
$Z$ and descend to (TDR.2).  Their images are finite because $P$
is finite.  This explicitly handles all possible projective
cocycles and does not assume that either factor has a linear
representation of $P$.

If a factor image had a proper invariant subspace, tensoring it
with the other entire factor would give a proper $J$-invariant
subspace of $V$.  If it had a nontrivial direct-sum system of
imprimitivity, tensoring each summand with the other factor would
give such a system for $J$.  Both are impossible.  This proves
primitivity of both images and also proves the final assertion
for any other given factorwise tensor decomposition.  A trivial
projective image on a space of dimension greater than one is
reducible, so it cannot be the image on $U$ when $\dim U>1$.
Finally $\dim U=d/\dim W\le d/2$, as required.
\end{proof}

\subsection*{D-TD12. Extraction and recursion through every layer}

\begin{theorem}[D-TD12]
Every finite primitive $P\subset\operatorname{PGL}(V)$, $d>1$,
has a factorization
\[
 V\simeq\bigotimes_{j=1}^s V_j,\qquad \dim V_j\ge2,
 \qquad s\le\lfloor\log_2 d\rfloor,
 \tag{TDR.6}
\]
with the following properties.
\begin{enumerate}
\item Each $V_j$ is either an irreducible represented projective
Heisenberg module, or a faithful irreducible module for a represented
finite quasisimple group.
\item The original group $P$ permutes all embedded full factor
algebras $\operatorname{End}(V_j)\otimes I$.  Whenever an operator
fixes one such algebra, its local operator normalizes the indicated
projective core.  For a quasisimple core it in fact normalizes the
represented perfect linear group.
\item The construction is a finite sequence of actual projective
quotients on smaller multiplicity spaces.  No multiplicity factor
of dimension greater than one is discarded or left outside (TDR.6).
\end{enumerate}
The pieces and the recursion need not be canonical.  Cores at a
later layer are represented in a projective quotient image; the
theorem does not assert that they lift as subgroups of the original
group $P$.
\end{theorem}

\begin{proof}
At a current primitive projective image $P_0$ choose a minimal
nontrivial normal subgroup $N$.  Such a subgroup is characteristically
simple: a characteristic subgroup of $N$ is normal in $P_0$.
We recall the elementary finite-group argument, rather than invoking
a classification of finite primitive groups.  Choose a minimal
nontrivial normal subgroup $S$ of $N$ and take all its distinct
automorphic images.  They are minimal normal subgroups of $N$;
distinct ones have trivial intersection and commute.  Their product
is characteristic and nontrivial, hence equals $N$.
If $S$ is abelian, this product is abelian; its characteristic
primary subgroups and power subgroups show that $N$ is an elementary
abelian $p$-group.  Otherwise each such $S$ has trivial center:
its center is normal in $N$ and minimality excludes a nontrivial
proper center.  The intersection of one factor with the product
of the others centralizes that factor, so it is trivial.  We
therefore have a direct product.  Every normal subgroup of one
factor is normal in this direct product, so minimality makes the
factor simple.  The factors are isomorphic because they were
obtained by automorphisms.  In this case
\[
                       N=S_1\times\cdots\times S_t
 \tag{TDR.7}
\]
with the $S_i$ nonabelian simple.  Conjugation by $P_0$ permutes
these factors transitively, since the product over one orbit
would otherwise contradict minimal normality of $N$.

Use D-TD11 for $N$, with finite special-linear lift $J$, scalar
kernel $Z$, inverse image $R$, and current space $W\otimes U$.
We now identify the extracted irreducible $R$-module $W$.

\paragraph{Elementary abelian case.}
Suppose $N$ is elementary abelian of exponent $p$.  Lifts of
$x,y\in N$ have scalar commutator, independent of their scalar
choices.  This defines an alternating biadditive pairing
\[
 \beta:N\times N\longrightarrow\mu_p,\qquad
 \beta(x,y)=[\widetilde x,\widetilde y].
 \tag{TDR.8}
\]
For example $\beta(x,y)^p=[\widetilde x^p,\widetilde y]=1$
because $\widetilde x^p\in Z$; commutator identities when
commutators are central prove biadditivity.  Conjugation by $J$
fixes the scalar values, so the radical is $P_0$-invariant.
Minimality of $N$ makes the radical either $0$ or all of $N$.
In the second case $R$ is abelian.  Its irreducible module $W$
has dimension one, contradicting D-TD11.  Thus $\beta$ is
nondegenerate.  A symplectic basis gives
\[
                        |N|=p^{2m},\qquad m\ge1.
 \tag{TDR.9}
\]

The projective image of $R$ on $W$ is faithfully $N$: an element
acting by a scalar on $W$ acts by a scalar on the full current
space, and so lies in $Z$.  Nondegeneracy also gives $Z(R)=Z$.
Every $x\in R\setminus Z$ has trace zero on $W$, since some
$y\in R$ conjugates it to $zx$ for a nonidentity scalar $z$.
Writing $a=\dim W$, character orthogonality now gives
\[
 1=\frac1{|R|}\sum_{x\in R}|\operatorname{tr}_W(x)|^2
   =\frac{|Z|a^2}{|R|},\qquad a=p^m.
 \tag{TDR.10}
\]
Hence this is precisely a represented irreducible projective
Heisenberg piece, including $p=2$.  Conjugation by every local
$A_g$ in (TDR.5) normalizes its actual represented linear lift,
because $R\triangleleft J$ and (TDR.3) identifies its matrices.
There is no use of a split symplectic normalizer extension.

\paragraph{Nonabelian case.}
Suppose (TDR.7) holds.  Let $A_i\subset R$ be the inverse image
of $S_i$ and put $E_i=A_i'$.  Since $A_i/Z$ is perfect,
\[
 A_i=E_iZ,\qquad
 E_i=[E_iZ,E_iZ]=[E_i,E_i].
 \tag{TDR.11}
\]
The group $E_i$ maps onto $S_i$, with central kernel $E_i\cap Z$.
Its center is exactly that kernel, so it is quasisimple.
For $i\ne j$ we have $[A_i,A_j]\subset Z$.  Fixing $y\in A_j$,
the map $x\mapsto[x,y]$ on $A_i$ is a homomorphism into the
abelian group $Z$.  It is trivial on $E_i=A_i'$.  Thus
$[E_i,A_j]=1$, and in particular all the $E_i$ commute.
Consequently
\[
              R=Z E_1\cdots E_t,
              \qquad R'=E_1\cdots E_t=:E.
 \tag{TDR.12}
\]
The intersections of components are central; their multiplication
map from $\prod E_i$ has central kernel.

The group $E$ acts irreducibly and faithfully on $W$, because
adjoining $Z$ only adjoins scalar matrices and $R$ is faithful
there.  Pull this representation back along $\prod E_i\to E$.
The product representation theorem
\cite[Section~3.2, Theorem~10 and proof, pp.~27--28]{OG:Serre1977}
gives
\[
                   W\simeq W_1\otimes\cdots\otimes W_t,
 \tag{TDR.13}
\]
where $W_i$ is an irreducible $E_i$-module.  Each is faithful:
an element acting trivially on $W_i$ acts trivially on $W$, hence
on the current full space, and is the identity.  Each has
dimension at least two, since a nontrivial perfect group has no
nontrivial one-dimensional representation.

An automorphism induced by $J$ permutes the simple factors of
$R/Z=N$, hence their inverse images $A_i$ and their derived
groups $E_i$.  The algebra generated by $E_i$ on $W$ is exactly
$\operatorname{End}(W_i)\otimes I$ by the same matrix-algebra
theorem.  Therefore every original local operator permutes these
embedded factor algebras.  If one is fixed, the corresponding
local operator conjugates the represented group $E_i$ to itself.
This proves the required normalization, including permutations,
without choosing incompatible component intertwiners.

\paragraph{Keeping all multiplicities.}
The coarse algebra $\operatorname{End}(W)\otimes I_U$ and its
commutant $I_W\otimes\operatorname{End}(U)$ are individually
normalized by the current image.  Retain the extracted Heisenberg
piece, or all the quasisimple pieces in (TDR.13).  If $\dim U=1$,
this layer is finished.  Otherwise D-TD11 gives a finite primitive
projective image on $U$ and a genuine quotient homomorphism onto
it, with strictly smaller dimension.  Apply exactly the same
construction to that quotient.

Every later factor algebra is embedded in a preceding commutant.
Its invariance or permutation under the smaller projective image
therefore pulls back to invariance or permutation under the entire
original group.  Scalar choices have no effect on conjugation.
More explicitly, a representative of an original operator on a
current multiplicity space differs by a scalar from an element
of the current finite special-linear lift; consequently it has
the same conjugation action on each newly represented core.
This proves normalization at later layers even though those
linear cores need not be subgroups of the original linear lift.

At every nonterminal step the multiplicity dimension is divided
by at least two.  Thus the process stops after at most
$\lfloor\log_2 d\rfloor$ extracted layers.  A terminal
multiplicity space cannot have dimension greater than one and
trivial projective image, by D-TD11; it has dimension one and
contributes only scalars.  Refining each nonabelian layer as in
(TDR.13), the product of all final factor dimensions is exactly
$d$, so their total number is also at most
$\lfloor\log_2 d\rfloor$.  This proves (TDR.6) and all claims.
\end{proof}

\subsection*{D-TD13. An unconditional whole-partner support interface}

\begin{corollary}[D-TD13]
Let $P\subset\operatorname{PGL}(U)$ be any finite primitive
projective image, $\dim U>1$, and perform D-TD12.  Let $m_*$ be
the largest integer $m\ge5$ for which a quasisimple factor at
\emph{any} recursion layer has simple quotient isomorphic to
$A_m$, with $m_*=1$ if no such factor occurs.  Using the local
quasisimple bounds of HD-FT20 and the Heisenberg bound D-TD08,
every nonscalar operator in this actual partner image satisfies
\[
 \frac{\nu_U(T)}{\dim U}\ge
      \frac1{\max\{80,m_*-1\}}.
 \tag{TDR.14}
\]
There is no factor count, factor degree, or unspecified residual
multiplicity hypothesis in this conclusion.  If a finite primitive
coefficient group acts factorwise on $V=W\otimes U$, its actual
projective image on $U$ satisfies these hypotheses.
\end{corollary}

\begin{proof}
D-TD12 gives exactly the complete normalized factor-algebra
decomposition required by HD-FT21, with no arbitrary bounded-degree
pieces.  Put $A=2$ in that theorem.  For clarity, the local floors
are $1/80$ for nonalternating quasisimple cores, $1/(m-1)$ for
$A_m$, $m\ge7$, $1/46$ for $A_5,A_6$, and $1/4$ for
Heisenberg cores.  If an operator fixes every factor algebra,
spectral convolution and a nonscalar local factor give their
minimum.  If it moves a factor cycle, the weighted tensor-cycle
bound of HD-FT18 gives at least $1/4$, independent of twists.
These statements yield (TDR.14).  The entire partner is covered:
the maximum $m_*$ includes all the recursively extracted layers,
not only the first normal core.  Primitivity inherited by the
partner is the last assertion of D-TD11.
\end{proof}

\paragraph{Actual quotient and source scope.}
If the coefficient group is given with its actual generating
product-one inertia tuple, applying any of the projective quotient
homomorphisms constructed above produces an actual quotient tuple.
Omitting its identity entries gives the connected quotient cover.
No assertion of twelve nonidentity entries in such a projection is
made.  At any original entry, spectral convolution transfers a
projective support bound on a factor to the actual moving rank of
the full coefficient operator, independently of its scalar phase.
A genuine scalar original entry has full moving rank.

The theorem does not assert individual quotient maps to quasisimple
components which are permuted by the original group; only their
whole coarse layer and the residual multiplicity have the stated
quotient maps.  It proves a complete complex representation
decomposition and the whole-partner support interface, not bounded
group orders, a bounded $m_*$, original coefficient-field descent,
polarization compatibility, or existence of any curve supplier.
In particular the earlier factorwise tensor-indecomposable
hypothesis of D-BI19 is not retroactively removed from that theorem.
